\appendix
\section{附录：核心代码与复现实验说明}

本附录说明项目代码结构、notebook 执行顺序和正文未展开的附录材料。完整代码位于项目目录中的 \texttt{notebooks/}、\texttt{src/} 和 \texttt{scripts/}，正文只展示必要片段，避免用大段代码替代方法说明。

\subsection{项目代码结构}

\begin{itemize}
  \item \texttt{src/config.py}：路径、随机种子、目标变量、泄漏列、聚类输入白名单和结果输出路径配置。
  \item \texttt{src/data\_utils.py}：数据下载、读取、数据集合规检查和摘要输出。
  \item \texttt{src/feature\_engineering.py}：工程特征生成、任务级特征排除和预处理管道构造。
  \item \texttt{src/model\_utils.py}：分类、回归、聚类评估与指标输出函数。
  \item \texttt{src/visualization.py}：EDA 图、模型评估图和聚类画像图绘制函数。
  \item \texttt{scripts/stage2\_5\_evidence\_enhancement.py}：预处理质量检查、PCA 解释方差和最终报告图表清单生成脚本。
\end{itemize}

\subsection{Notebook 执行顺序}

复现实验时按以下顺序执行 notebook：

\begin{enumerate}
  \item \texttt{notebooks/00\_dataset\_selection\_and\_compliance.ipynb}
  \item \texttt{notebooks/01\_data\_preprocessing\_and\_eda.ipynb}
  \item \texttt{notebooks/02\_classification\_high\_risk.ipynb}
  \item \texttt{notebooks/03\_regression\_productivity.ipynb}
  \item \texttt{notebooks/04\_clustering\_lifestyle\_profiles.ipynb}
  \item \texttt{notebooks/05\_result\_summary\_for\_report.ipynb}
\end{enumerate}

\subsection{附录图表}

图 \ref{fig:appendix-numeric-histograms} 展示主要数值变量直方图，作为第 4 章 EDA 的补充。图 \ref{fig:appendix-behavior-scatter} 展示行为变量与结果变量的散点或分箱关系，用于辅助观察变量关系。

\maybefigure{eda_numeric_histograms.png}{主要数值变量分布直方图}{fig:appendix-numeric-histograms}

\maybefigure{eda_behavior_outcome_scatter.png}{行为变量与结果变量关系图}{fig:appendix-behavior-scatter}

图 \ref{fig:appendix-productivity-prediction} 展示 \texttt{productivity\_score} 的真实值与预测值对比。该图作为生产力弱预测结果的补充材料。图 \ref{fig:appendix-classification-importance} 展示分类任务置换重要性，用于补充解释分类模型主要依赖的输入变量。

\maybefigure{regression_productivity_observed_vs_predicted.png}{生产力得分真实值与预测值对比}{fig:appendix-productivity-prediction}

\maybefigure{classification_permutation_importance.png}{分类任务置换重要性}{fig:appendix-classification-importance}

\subsection{附录表格与结果文件}

PCA 解释方差表作为附录材料列出前 10 个主成分，完整结果见 \texttt{results/pca\_explained\_variance.csv}。

\begin{table}[H]
\centering
\caption{PCA 解释方差附录表（前 10 个主成分）}
\label{tab:pca-explained-variance}
\begin{tabular}{rrr}
\toprule
主成分 & 解释方差比例 & 累计解释方差 \\
\midrule
PC1 & 0.2694 & 0.2694 \\
PC2 & 0.1547 & 0.4241 \\
PC3 & 0.1366 & 0.5607 \\
PC4 & 0.1258 & 0.6865 \\
PC5 & 0.1030 & 0.7895 \\
PC6 & 0.0759 & 0.8654 \\
PC7 & 0.0430 & 0.9084 \\
PC8 & 0.0425 & 0.9509 \\
PC9 & 0.0274 & 0.9783 \\
PC10 & 0.0152 & 0.9935 \\
\bottomrule
\end{tabular}
\end{table}


分类调参完整结果保存在 \texttt{results/classification\_cv\_results.csv}，数字依赖回归调参完整结果保存在 \texttt{results/regression\_digital\_dependence\_cv\_results.csv}。这些 CSV 文件列数较多，正文不直接排版完整表格，避免影响阅读；正式提交时可作为附录电子材料与 notebook 一并提供。

\subsection{参考资料与数据来源}

\begin{enumerate}
  \item Digital Lifestyle Benchmark Dataset (2025).
  \item scikit-learn 官方文档。
  \item 《大数据分析与应用》课程考查报告要求。
\end{enumerate}
